
This thesis investigated how an AI-generated structured DeFi intent can
be verified against authoritative execution relationships before blockchain
submission. The central problem was that individually valid protocols,
contracts, functions, and assets can still form an inconsistent execution
structure. Neither a well-formed model response nor wallet authorization
establishes that these entities belong together for the proposed operation.
The objective was therefore to investigate a maintained, domain-specific
property graph as a pre-execution verification mechanism and to connect
its decisions to controlled ERC-4337 wallet execution.

The developed concept separates natural-language intent generation,
deterministic validation, verification-knowledge acquisition,
action-specific relationship verification, and controlled execution.
The language model produces a candidate intent, while execution-specific
contracts, functions, and protocol resources are resolved from admitted
graph knowledge. The prototype implements this approach using Neo4j for
selected Uniswap V3 single-hop swaps and Aave V3 supply, withdrawal,
borrowing, and repayment operations on Ethereum Sepolia. Approved
operations are retained server-side and reverified before transaction
preparation.

\medskip
\noindent\textbf{Contributions and Main Findings}

The principal contribution is an action-specific verification architecture
that evaluates the connections among execution entities within the
applicable network and protocol deployment. The graph serves as maintained
verification knowledge against which the structured operation is checked.
This makes the relationships required for approval explicit and separates
their verification from probabilistic language interpretation. The approach
also distinguishes evidence of an unsupported operation from insufficient
evidence, preventing either condition from authorizing execution.

The controlled relationship experiment provides the clearest evidence
for this contribution. Both the full verifier and an entity-existence
baseline approved 12 valid controls. After individual relationship
deletions, the full verifier rejected all 60 altered paths, whereas the
baseline approved all 60. Since the entities remained present, the
comparison demonstrates the importance of checking their required
connections. It supports the tested verification patterns without
establishing that a graph database is the only suitable implementation.

The evaluation also examined graph verification independently of
language-model output. All 340 direct reference-intent calls, covering
112 distinct graph inputs, agreed with the decisions and selected
execution evidence derived from the admitted snapshot. Of these calls,
338 completed within two seconds, corresponding to 99.41\%. The first
occurrences of distinct inputs met the threshold in 110 of 112 cases,
or 98.21\%. These results satisfy the prototype-specific R6 target for
the recorded workload. Separately, 36 knowledge-maintenance observations
across 12 controlled conditions matched the predefined expectations,
including non-approval under stale or incomplete evidence.

The balanced natural-language study produced 534 final-status matches
across 600 prompts, corresponding to 89.00\% agreement with the draft
references. The more restrictive joint-reference measure matched in
529 cases, or 88.17\%, while four attempts ended in provider quota errors.
These findings show that the complete pipeline remains less reliable
than the graph checks considered separately. The research objective
has therefore been achieved for relationship verification within the
admitted and evaluated scope, while integration with controlled wallet
execution was initially demonstrated at the application handover level.
A separate November regression subsequently demonstrated exact-batch,
one-use authorization through actual local Kernel and EntryPoint contracts.

\medskip
\noindent\textbf{Limitations}

The most important limitation concerns preservation of the original
request. In one evaluated case, the language model omitted an explicit
intermediate-token condition and produced a direct-swap intent. The
graph verifier approved that reduced operation because its relationships
were supported by the admitted knowledge. No transaction was submitted.
This demonstrates that relationship correctness does not establish
intent correctness: a verifier can correctly assess the operation it
receives while a user constraint has already been lost. The evaluation also has methodological boundaries. The natural-language
references were assistant-drafted without documented independent
item-level review, and the joint measure checks only selected aspects
of interpretation. The graph references were derived from the admitted
snapshot rather than an independent validation of every protocol fact.
Prompts, graph inputs, and mutation fixtures contain dependent
observations. Furthermore, the earlier 400-prompt study and the
600-prompt follow-up used different datasets and do not establish
causal improvement.

Verification remains dependent on the quality and maintenance of its
reference knowledge. Synchronization procedures automate acquisition
when invoked, but continuous source monitoring is not implemented.
Manually maintained function mappings, non-atomic updates, incomplete
removal of obsolete relationships, and limited representation of live
reserve restrictions constrain the meaning of approval. The graph
establishes consistency with admitted knowledge rather than complete
current transaction feasibility. Traceability and deployment assurance
also remain incomplete. The original 28 wallet-policy and handover tests
included simulated responses. The later local authorization regression
supports R4 for exact approved batches and one-use restrictions, but does
not establish Sepolia installation, public-bundler acceptance, cumulative
budgets across approvals, or successful testnet settlement.
The subsequent traceability regression correctly localized all 60
controlled relationship deletions and reported their expected and
observed conditions. However, approved controls still lacked an explicit
reason or verification-rule identifier, and some broader graph results
remained unlocalized. The improvement therefore supports traceability
for the tested mutation cases rather than complete R5 coverage across
all outcomes.

\medskip
\noindent\textbf{Future Work}

An immediate research direction is stronger preservation of explicit
user constraints throughout interpretation, validation, and transaction
preparation. This should be evaluated using independently reviewed
references and requests kept separate from development. Evaluation
should also distinguish failure to extract a constraint from failure
to enforce a correctly extracted constraint. Further work should strengthen knowledge maintenance through consistent
updates, dependency-aware freshness checks, and validation of relevant
protocol state. Extending the approach to additional protocols,
networks, and actions would help determine which relationship patterns
are reusable and which require protocol-specific verification rules.
Beyond the initial scope, multi-step operations offer a further research
question: how can individually valid actions be verified as a consistent
sequence when the output and state changes of one action affect the
requirements of the next?

Further work should independently audit and deploy the locally tested
authorization mechanism, validate public-bundler compatibility, and add
authenticated browser-account binding, cumulative budgets across approvals,
and persistent receipt reconciliation.
Comparing the property-graph implementation with equally expressive
rule-based or relational alternatives would also clarify the practical
benefits and costs of the chosen representation.

The thesis shows that explicit, action-specific relationship verification
can provide a useful safeguard between AI-generated DeFi intents and
blockchain execution. Its contribution is a verifiable connection between
an interpreted operation and admitted execution knowledge. The security
value of that connection depends on preserving the user's constraints,
maintaining the underlying knowledge, and enforcing the approved
operation at the execution boundary.